% ===========================================================================
\chapter{Modelagem das Personas Arquetípicas}
\label{cap_personas_arquetipicas}
% ===========================================================================

A modelagem de personas para o SIGEA-GTT fundamenta-se em investigação qualitativa conduzida em cooperação multidisciplinar entre o Departamento de Computação (DCOMP) e o corpo docente e discente do Departamento de Enfermagem da Universidade Federal de Sergipe (UFS). Foram realizadas entrevistas semiestruturadas com docentes de disciplinas de Gestão em Enfermagem e Segurança do Paciente, além de observações diretas de sessões de ensino prático de auditoria em ambiente laboratorial com graduandos.

A partir da triangulação desses dados, emergiram padrões consistentes de comportamento, modelos mentais, objetivos pedagógicos e barreiras emocionais, sintetizados em três personas complementares e uma anti-persona delimitadora.

% ===========================================================================
\section{Persona Primária: Beatriz Matos (Discente de Enfermagem)}
\label{sec_persona_beatriz}
% ===========================================================================

\begin{figure}[!ht]
\centering
\begin{tikzpicture}
    \node[draw=clinical-700, fill=clinical-100!30, rounded corners=6pt, inner sep=12pt, line width=1pt, text width=0.92\textwidth] (card) {
        \begin{minipage}{0.28\textwidth}
            \centering
            \begin{tikzpicture}
                \draw[fill=clinical-700!15, draw=clinical-700, line width=1.5pt] (0,0) circle (1.3cm);
                \node at (0,0.3) {\scalebox{2.8}{\color{clinical-700}$\mathbf{\Omega}$}};
                \node at (0,-0.6) {\scriptsize\bfseries BEATRIZ MATOS};
            \end{tikzpicture}
            \vspace{4pt}
            
            \footnotesize
            \textbf{Idade}: 22 anos \\
            \textbf{Curso}: Enfermagem -- UFS \\
            \textbf{Período}: 7º Período \\
            \textbf{Perfil RBAC}: \texttt{STUDENT} \\
            \textbf{Dispositivo}: Desktop Laboratório
        \end{minipage}
        \hfill
        \begin{minipage}{0.67\textwidth}
            \footnotesize
            \textbf{``Quero aprender a rastrear eventos adversos em um ambiente seguro, onde errar seja uma oportunidade de aprendizado sobre causas sistêmicas, e não motivo de medo ou punição.''}
            \vspace{6pt}
            \hrule
            \vspace{6pt}
            \textbf{Metas e Objetivos Principais}:
            \begin{itemize}\setlength{\itemsep}{1pt}
                \item Dominar a metodologia IHI-GTT antes de ingressar no internato hospitalar;
                \item Cumprir a auditoria retrospectiva no limite estrito de 20 minutos sem pânico;
                \item Rastrear gatilhos ocultos em evoluções, prescrições e laudos laboratoriais;
                \item Aplicar ferramentas da qualidade (Ishikawa 6M, GUT, 5W2H, PDCA) de forma conectada ao caso clínico.
            \end{itemize}
            \textbf{Principais Dores e Frustrações}:
            \begin{itemize}\setlength{\itemsep}{1pt}
                \item Bloqueio cognitivo frente a prontuários físicos volumosos e desorganizados;
                \item Ansiedade intensa com contagem de tempo sem avisos graduais;
                \item Medo arraigado da cultura punitiva de culpar o profissional individual;
                \item Falta de retorno imediato sobre quais gatilhos deixou de detectar.
            \end{itemize}
        \end{minipage}
    };
\end{tikzpicture}
\caption{Ficha Sintética da Persona Primária: Beatriz Matos}
\label{fig_persona_beatriz}
\fonte{Elaborado pelo autor (2026).}
\end{figure}

\subsection{Perfil Sociodemográfico e Contexto Acadêmico}
Beatriz tem 22 anos e cursa o 7º período da graduação em Enfermagem no campus de São Cristóvão da UFS. Divide sua semana entre aulas teóricas pela manhã e estágios práticos supervisionados nas enfermarias do Hospital Universitário de Sergipe (HU-UFS) durante a tarde. Costuma estudar nos laboratórios de informática do campus em computadores com monitores HD ($1280 \times 720$) e Full HD ($1920 \times 1080$), intercalando leitura de artigos científicos e preparação de planos de cuidados.

Possui fluência digital com redes sociais, buscadores acadêmicos e aplicativos de produtividade pessoal, mas ressente-se da lentidão e complexidade confusa dos sistemas hospitalares legados aos quais é exposta durante os plantões de estágio.

\subsection{Rotina e Modelo Mental de Auditoria}
Ao receber uma tarefa de auditoria clínica, Beatriz costumava imprimir até trinta páginas de prontuários simulados. Sua rotina envolvia destacar frases com canetas marca-texto fluorescentes enquanto controlava o tempo com o cronômetro do próprio smartphone. Esse modelo analógico gerava frequentes interrupções: o alarme do celular disparava de forma abrupta e ruidosa, interrompendo sua linha de raciocínio, e ao final da tarefa suas folhas estavam repletas de anotações dispersas que não se comunicavam com as matrizes de Ishikawa ou 5W2H distribuídas em folhas separadas.

Beatriz anseia por um modelo mental integrado: visualização do prontuário em painéis tabulados estruturados (admissão, evoluções diárias, prescrições, exames laboratoriais numéricos com faixas de referência e cirurgias), seleção de texto direta com vinculação imediata a um gatilho canônico do IHI, e preenchimento progressivo das ferramentas da qualidade na mesma tela unificada.

\subsection{Citação Característica}
\begin{quote}
\textit{``Nos primeiros dias de estágio no hospital, vi um paciente receber um medicamento com dosagem incorreta porque a prescrição estava ilegível. Fiquei paralisada. Quero praticar muito no SIGEA-GTT para que meus olhos aprendam a enxergar esses gatilhos instantaneamente, antes que o dano atinja o paciente real.''}
\end{quote}

% ===========================================================================
\section{Persona Secundária: Prof. Dr. Ricardo Sobral (Docente Titular)}
\label{sec_persona_ricardo}
% ===========================================================================

\begin{figure}[!ht]
\centering
\begin{tikzpicture}
    \node[draw=hospital-green, fill=hospital-green!10, rounded corners=6pt, inner sep=12pt, line width=1pt, text width=0.92\textwidth] (card) {
        \begin{minipage}{0.28\textwidth}
            \centering
            \begin{tikzpicture}
                \draw[fill=hospital-green!20, draw=hospital-green, line width=1.5pt] (0,0) circle (1.3cm);
                \node at (0,0.3) {\scalebox{2.8}{\color{hospital-green}$\mathbf{\Psi}$}};
                \node at (0,-0.6) {\scriptsize\bfseries PROF. DR. RICARDO};
            \end{tikzpicture}
            \vspace{4pt}
            
            \footnotesize
            \textbf{Idade}: 48 anos \\
            \textbf{Cargo}: Docente Adjunto \\
            \textbf{Titulação}: Doutor em Saúde \\
            \textbf{Perfil RBAC}: \texttt{PROFESSOR} \\
            \textbf{Dispositivo}: Workstation / Ultrawide
        \end{minipage}
        \hfill
        \begin{minipage}{0.67\textwidth}
            \footnotesize
            \textbf{``Corrigir manualmente cinquenta auditorias em papel consome minhas noites. Preciso de uma ferramenta que compare as entregas dos alunos com o meu gabarito e consolide indicadores científicos sem erros de fórmulas.''}
            \vspace{6pt}
            \hrule
            \vspace{6pt}
            \textbf{Metas e Objetivos Principais}:
            \begin{itemize}\setlength{\itemsep}{1pt}
                \item Criar turmas padronizadas (\textit{``Disciplina - Turma - Período''});
                \item Publicar prontuários fictícios de alta fidelidade com gatilhos calibrados;
                \item Corrigir com espelho comparativo (Sensibilidade, Falsos Positivos);
                \item Atribuir nota de 0 a 10 com parecer formativo obrigatório;
                \item Consolidar automaticamente taxas do IHI ($TI_{1000}$, $TI_{100}$, $P_{EA}$).
            \end{itemize}
            \textbf{Principais Dores e Frustrações}:
            \begin{itemize}\setlength{\itemsep}{1pt}
                \item Sobrecarga estafante na conferência manual de submissões discentes;
                \item Dificuldade de manter controle de versões de planilhas de notas no Excel;
                \item Insegurança quanto à colisão ou duplicidade de turmas no mesmo período letivo;
                \item Alunos que entregam auditorias incompletas por falta de acompanhamento de prazos.
            \end{itemize}
        \end{minipage}
    };
\end{tikzpicture}
\caption{Ficha Sintética da Persona Secundária: Prof. Dr. Ricardo Sobral}
\label{fig_persona_ricardo}
\fonte{Elaborado pelo autor (2026).}
\end{figure}

\subsection{Perfil Sociodemográfico e Trajetória Acadêmica}
O Prof. Dr. Ricardo Sobral tem 48 anos, é enfermeiro de formação com mestrado e doutorado em Saúde Coletiva pela Fundação Oswaldo Cruz (FIOCRUZ). Atua há dezesseis anos como professor titular das disciplinas de Gestão em Enfermagem e Segurança do Paciente na UFS, além de liderar o Núcleo de Segurança do Paciente (NSP) do HU-UFS. Em sua sala de trabalho, utiliza uma estação desktop com monitor \textit{Ultrawide} ($2560 \times 1080$), permitindo abrir simultaneamente a literatura científica, prontuários de exemplo e painéis avaliativos.

Ricardo é um docente extremamente rigoroso quanto à metodologia científica. Conhece a fundo os manuais do IHI e as categorias de dano do NCC MERP \cite{nccmerp2001}, defendendo que os indicadores hospitalares de eventos adversos devem ser ensinados com precisão matemática.

\subsection{Gargalos Docentes e Necessidades de Design}
A maior fonte de exaustão profissional de Ricardo reside na etapa avaliativa semestral. Com turmas de até quarenta discentes, cada atividade prática resulta em quarenta relatórios impressos de dez a quinze páginas. A conferência manual de quais dos 53 gatilhos canônicos foram acertados por cada estudante, a contagem de falsos positivos e a verificação do cálculo do score da Matriz GUT ($G \times U \times T$) consomem dezenas de horas de trabalho repetitivo.

Ricardo necessita de uma plataforma que:
\begin{enumerate}
    \item Impeça a criação desordenada de turmas, exigindo a máscara canônica \texttt{"Disciplina - Turma - Período"} e vinculando-o obrigatoriamente como docente titular;
    \item Ofereça um painel de correção lado a lado: as respostas do discente justapostas ao gabarito do professor, com cálculo instantâneo da concordância de classificação de dano;
    \item Bloqueie o encerramento do período letivo caso existam atividades com notas pendentes;
    \item Calcule e plote gráficos com as taxas oficiais do IHI: Taxa de Eventos Adversos por 1.000 Pacientes-Dia ($TI_{1000}$), Taxa por 100 Admissões ($TI_{100}$) e Prevalência de Pacientes com Dano ($P_{EA}$).
\end{enumerate}

\subsection{Citação Característica}
\begin{quote}
\textit{``O ensino de auditoria clínica precisa abandonar o artesanato de planilhas. Quero gastar o meu tempo orientando os alunos sobre como propor melhorias estruturais no ciclo PDCA, e não somando notas de papel com caneta vermelha.''}
\end{quote}

% ===========================================================================
\section{Persona Terciária: Carlos Eduardo Santos (Administrador de TI)}
\label{sec_persona_carlos}
% ===========================================================================

\begin{figure}[!ht]
\centering
\begin{tikzpicture}
    \node[draw=clinical-500, fill=clinical-500!10, rounded corners=6pt, inner sep=12pt, line width=1pt, text width=0.92\textwidth] (card) {
        \begin{minipage}{0.28\textwidth}
            \centering
            \begin{tikzpicture}
                \draw[fill=clinical-500!20, draw=clinical-500, line width=1.5pt] (0,0) circle (1.3cm);
                \node at (0,0.3) {\scalebox{2.8}{\color{clinical-500}$\mathbf{\Delta}$}};
                \node at (0,-0.6) {\scriptsize\bfseries CARLOS SANTOS};
            \end{tikzpicture}
            \vspace{4pt}
            
            \footnotesize
            \textbf{Idade}: 35 anos \\
            \textbf{Cargo}: Analista de TI \\
            \textbf{Lotação}: DCOMP / UFS \\
            \textbf{Perfil RBAC}: \texttt{ADMIN} \\
            \textbf{Dispositivo}: Estação WUXGA
        \end{minipage}
        \hfill
        \begin{minipage}{0.67\textwidth}
            \footnotesize
            \textbf{``Minha prioridade é garantir que o sistema seja inexpugnável, que apenas a comunidade acadêmica legítima acerte o login e que os dados simulados estejam 100\% em conformidade com a LGPD.''}
            \vspace{6pt}
            \hrule
            \vspace{6pt}
            \textbf{Metas e Objetivos Principais}:
            \begin{itemize}\setlength{\itemsep}{1pt}
                \item Garantir validação estrita de e-mails com sufixo \texttt{@academico.ufs.br};
                \item Segregar perfis de forma estanque (\texttt{ADMIN}, \texttt{PROFESSOR}, \texttt{STUDENT});
                \item Forçar redefinição de senha no primeiro login de novos usuários;
                \item Manter o catálogo dos 6 módulos e 53 gatilhos canônicos atualizado;
                \item Garantir isolamento absoluto contra prontuários reais de pacientes do SUS.
            \end{itemize}
            \textbf{Principais Dores e Frustrações}:
            \begin{itemize}\setlength{\itemsep}{1pt}
                \item Sobrecarga de chamados de suporte por senhas esquecidas;
                \item Inconsistências de cadastro (ex.: docentes com campo matrícula preenchido indevidamente);
                \item Falhas de conformidade com a Lei Geral de Proteção de Dados (LGPD).
            \end{itemize}
        \end{minipage}
    };
\end{tikzpicture}
\caption{Ficha Sintética da Persona Terciária: Carlos Eduardo Santos}
\label{fig_persona_carlos}
\fonte{Elaborado pelo autor (2026).}
\end{figure}

\subsection{Perfil e Responsabilidades Técnicas}
Carlos Eduardo tem 35 anos e atua como Analista de Tecnologia da Informação do DCOMP-UFS, sendo responsável pela infraestrutura dos servidores de aplicação, laboratórios acadêmicos e bancos de dados acadêmicos. Trabalha em uma estação corporativa com monitor WUXGA ($1920 \times 1200$) em ambiente Linux.

Carlos zela pela segurança da informação, integridade referencial dos dados e obediência às diretrizes centrais de tecnologia da Universidade. Para ele, uma interface administrativa eficiente deve ser objetiva, livre de redundâncias e com validações automáticas que impeçam a inserção de dados inconsistentes antes mesmo de atingirem a camada de persistência.

\subsection{Citação Característica}
\begin{quote}
\textit{``Se o sistema rejeitar qualquer e-mail que não seja @academico.ufs.br e impedir que professores recebam matrícula de aluno por engano, eliminamos 90\% dos chamados de suporte técnico do DCOMP.''}
\end{quote}

% ===========================================================================
\section{Anti-Persona: O Auditor Punitivo Tradicional}
\label{sec_anti_persona}
% ===========================================================================

Para blindar o SIGEA-GTT contra distorções funcionais e desvios de finalidade, foi modelada a anti-persona denominada \textbf{O Auditor Punitivo Tradicional}. Este perfil sintetiza a figura do fiscal burocrático hospitalar que utiliza a auditoria de prontuários com o propósito persecutório de identificar profissionais de enfermagem ou medicina para aplicação de sanções disciplinares, cortes salariais ou advertências formais.

O SIGEA-GTT \textbf{não} foi projetado para esse perfil pelos seguintes motivos fundamentais:
\begin{itemize}
    \item \textbf{Cultura Justa de Segurança}: Conforme preceitos do IHI e da Organização Mundial da Saúde (OMS), a segurança do paciente fundamenta-se na identificação de falhas latentes nos processos hospitalares (comunicação ineficaz, sobrecarga de leitos, calibração de bombas de infusão), e não na culpabilização individual do operador na ponta assistencial;
    \item \textbf{Natureza Estritamente Educacional}: A plataforma opera com dados sintéticos e fictícios de alta fidelidade persistidos em colunas semiestruturadas \texttt{JSONB}, sendo terminantemente desvinculada das bases de dados de produção assistencial do HU-UFS ou do Ministério da Saúde;
    \item \textbf{Foco na Aprendizagem Formativa}: O módulo avaliativo docente prioriza o parecer pedagógico e a compreensão dos fatores causais no Diagrama de Ishikawa 6M e Matriz GUT, repudiando qualquer métrica de ranqueamento punitivo.
\end{itemize}
